\documentclass[10pt]{article}
\linespread{1.05}

\usepackage{cmap}      % makes the PDF text copy/paste- and search-able
\usepackage[T1]{fontenc}
\usepackage[utf8]{inputenc}
\usepackage{amsfonts,amssymb,amsmath}
\usepackage[lmargin=1in,rmargin=1in,tmargin=0.9in,bmargin=0.9in]{geometry}
\usepackage{enumitem}
\usepackage[hidelinks]{hyperref}

\pagestyle{empty}
\setlength{\parindent}{0pt}
\sloppy % prefer loose spacing over lines running into the margin

% --- helper macros -----------------------------------------------------
% A section heading preceded by a full-width rule.
\newcommand{\cvsection}[1]{%
  \par\vspace{0.24cm}%
  \noindent\rule{\textwidth}{0.4pt}%
  \par\vspace{0.12cm}%
  \noindent\textbf{\large #1}%
  \par\vspace{0.12cm}%
}
% A bulleted list with tight spacing, used for every section body.
\newenvironment{cvlist}%
  {\begin{list}{$\bullet$}{%
     \setlength{\leftmargin}{1.1em}\setlength{\labelwidth}{0.8em}%
     \setlength{\itemsep}{0.12cm}\setlength{\parsep}{0pt}%
     \setlength{\topsep}{0pt}\setlength{\partopsep}{0pt}}}%
  {\end{list}}
% -----------------------------------------------------------------------

\begin{document}

%% ======================= HEADER =======================
\begin{minipage}[t]{0.48\textwidth}
\vspace{0pt}
{\Large\textbf{Dominic Bunnett}}\\[0.30cm]
\href{https://dombunnett.github.io}{\em dombunnett.github.io}\\
ORCID: \href{https://orcid.org/0009-0007-6921-3177}{\em 0009-0007-6921-3177}
\end{minipage}\hfill
\begin{minipage}[t]{0.48\textwidth}
\vspace{0pt}
Sekretariat MA 3-2\\
Institut f\"ur Mathematik\\
Technische Universit\"at Berlin\\
Stra\ss{}e des 17. Juni 136\\
10623 Berlin, Germany\\
Email: {\em bunnett@math.tu-berlin.de}
\end{minipage}

\vspace{0.35cm}

%% ======================= RESEARCH INTERESTS =======================
\cvsection{Research interests}
Moduli theory and geometric invariant theory, at the interface of algebraic and
discrete geometry, with an emphasis on computational and algorithmic problems.

%% ======================= EDUCATION =======================
\cvsection{Education}
\begin{cvlist}
\item {\bf Freie Universit\"at Berlin} \quad October 2015 -- July 2019\\
{\em Doctor of Philosophy in Mathematics}\\
Thesis: \href{https://refubium.fu-berlin.de/handle/fub188/27333}{Moduli of hypersurfaces in toric orbifolds}\\
Advisor: Prof.\ Victoria Hoskins
\item {\bf University of Warwick} \quad 2011 -- 2015\\
{\em BSc and MMath}, $1^{\text{st}}$ class\\
Thesis: Hodge theory and the Torelli theorem for curves\\
Advisor: Prof.\ Miles Reid
\end{cvlist}

%% ======================= POSITIONS =======================
\cvsection{Positions held}
\begin{cvlist}
\item {\bf Technische Universit\"at Berlin} \quad October 2021 -- present\\
Postdoc, research group Algorithmic Algebra\\
Mentor: Peter B\"urgisser
\item {\bf The Open University, UK} \quad September 2024 -- present\\
Associate Lecturer, {\em Applied Computational Mathematics}
% TODO: add the OU module code (e.g. MST3xx) if you want it, and confirm
% whether the contract end in October 2026 should be shown as an end date.
\item {\bf Makerere University, Kampala} \quad April 2022 -- January 2023\\
Lecturer
\item {\bf Technische Universit\"at Berlin} \quad August 2019 -- May 2021\\
Postdoc within the Einstein Visiting Fellowship of Bernd Sturmfels
(Einstein Foundation Berlin); initially as a Stipendiat, subsequently employed\\
Mentors: Michael Joswig and Bernd Sturmfels
\item {\bf University of Oxford} \quad April -- June 2019\\
Visiting Researcher, hosted by Frances Kirwan
\item {\bf Freie Universit\"at Berlin} \quad September 2015 -- March 2019\\
Doktorand / Wissenschaftlicher Mitarbeiter\\
Advisor: Victoria Hoskins
\end{cvlist}

%% ======================= PUBLICATIONS =======================
\cvsection{Publications}
\enlargethispage{-2\baselineskip} % keep the last entry off the page break
\begin{cvlist}
\item On the moduli of hypersurfaces in toric orbifolds.
{\em Proceedings of the Edinburgh Mathematical Society}, 67(2):577--616, 2024.
\href{https://doi.org/10.1017/S0013091524000166}{doi:10.1017/S0013091524000166},
\href{https://arxiv.org/abs/1906.00272}{arXiv:1906.00272}
\item Generalised cone complexes and tropical moduli in polymake.
{\em Joint with Michael Joswig and Julian Pfeifle.}
In {\em Proceedings of the 2023 International Symposium on Symbolic and
Algebraic Computation} (ISSAC '23), pages 100--106. ACM, New York, 2023.
\href{https://dl.acm.org/doi/abs/10.1145/3597066.3597093}{doi:10.1145/3597066.3597093}
\item Determinantal representations of the cubic discriminant.
{\em Joint with Hanieh Keneshlou.}
{\em Le Matematiche}, 75(2):489--505, 2020.
\href{https://lematematiche.dmi.unict.it/index.php/lematematiche/article/view/1997}{journal},
\href{https://arxiv.org/abs/1909.05579}{arXiv:1909.05579}
\end{cvlist}

%% ======================= PREPRINTS =======================
\cvsection{Preprints}
\begin{cvlist}
\item Embeddings of weighted projective spaces.
{\em Joint with Praise Adeyemo and Fabi\'an Levic\'an-Santib\'a\~nez.}
\href{https://arxiv.org/abs/2510.05076}{arXiv:2510.05076} (submitted)
\item Moduli of Bridgeland semistable holomorphic triples.
{\em Joint with Alejandra Rinc\'on-Hidalgo.}
\href{https://arxiv.org/abs/2102.04995}{arXiv:2102.04995} (submitted)
\end{cvlist}

%% ======================= IN PREPARATION =======================
\cvsection{In preparation}
\begin{cvlist}
\item Finiteness of automorphisms via Bott vanishing.
{\em Joint with Caroline Namanya.} (to appear on the arXiv, September 2026)
\item K-moduli of exceptional del Pezzo hypersurfaces.
{\em Joint with Josh Jackson, Jes\'us Mart\'inez-Garc\'ia and Charlotte Satchwell.}
\item Geometry of moduli stacks of hypersurfaces. {\em Joint with Josh Jackson.}
\item Hilbert--Mumford criteria for hypersurfaces.
\end{cvlist}

%% ======================= AWARDS =======================
\cvsection{Awards and competitive research programmes}
\begin{cvlist}
\item {\bf Research in Teams}, Erwin Schr\"odinger International Institute for
Mathematics and Physics (ESI), Vienna, May -- June 2025.\\
Project: {\em Projective geometry, combinatorics and representation theory from
polytopes}, with Praise Adeyemo, Livia Campo and Bal\'azs Szendr\H{o}i.
\item {\bf Research in Pairs}, Mathematisches Forschungsinstitut Oberwolfach,
2021.\\
Awarded jointly with Alejandra Rinc\'on-Hidalgo.
% TODO: confirm exact dates (CV said Sept 2021, your old TU page said June 2021)
% and add the official project title from MFO Annual Report 2021, p.98.
\item {\bf Clay Mathematics Institute}, full funding for the workshop
{\em Equivariant methods in geometry}, University of Ibadan, 2024.
\item {\bf MATH+ travel grants}, for research and teaching visits to Makerere
University and AIMS Ghana.
\item {\bf International Centre for Theoretical Physics}, funding for research
visits to Trieste (2019, 2022) and for the EAUMP--ICTP summer school,
University of Dar-es-Salaam, 2018.
\end{cvlist}

%% ======================= TALKS =======================
\cvsection{Invited talks}
\begin{cvlist}
\item \href{https://www.newton.ac.uk/event/ggiw02/}{AGGITatE 2026}, workshop of the
Isaac Newton Institute satellite programme {\em Algebraic Groups, Geometry,
Invariants and Related Topics}, University of Essex, July 2026.\\
\emph{Finite automorphisms of toric orbifold hypersurfaces.}
\item \emph{Embeddings of weighted projective spaces and rectangular simplices.}
Discrete Mathematics / Geometry seminar, TU Berlin, November 2025.
\item \emph{Embedding weighted projective spaces.}
\href{https://homepage.univie.ac.at/balazs.szendroi/?page_id=783}{Workshop in Algebraic Geometry},
Erwin Schr\"odinger International Institute for Mathematics and Physics,
Vienna, 2025.
% TODO: month, if you want it (the ESI stay ran 17 May - 25 June 2025)
\item \emph{Families of hypersurfaces from polytopes.} Conference
\href{https://icms.ac.uk/activities/workshop/algebra-and-geometry-from-africa/}{\em Algebra and Geometry from Africa},
ICMS, Edinburgh, May 2025.
\item \emph{The Frobenius problem, simplices and moduli spaces.}
Mathematics and Statistics Colloquium, The Open University, February 2025.
\item \href{https://www.youtube.com/watch?v=7kdnNRKsQLM}{\em Moduli of coverings branched along a hypersurface.}
Programme seminar, {\em New Equivariant Methods in Algebraic and Differential
Geometry}, Isaac Newton Institute, Cambridge, July 2024.
\item \emph{Building collaborations with Africa.} Symposium
\href{https://ias.uva.nl/content/events/2024/06/africa-symposium.html}{\em Building partnerships with institutions and researchers in Africa},
Institute for Advanced Study, University of Amsterdam, June 2024.
\item \emph{Geometry of some moduli spaces of hypersurfaces.}
\href{https://cow.alggeo.xyz/cow2023}{COW Seminar}, University of Bath, March 2024.
\item \emph{Moduli spaces of hypersurfaces and their topology.}
\href{https://researchseminars.org/talk/EssexMaths/97/}{Mathematics Essex Seminar Series},
University of Essex, December 2023.
\item \emph{Stacky fans and tropical moduli in polymake.}
\href{https://www.issac-conference.org/2023/program.php}{ISSAC 2023},
Troms\o, Norway, July 2023.
\item \emph{Moduli spaces of weighted hypersurfaces.}
\href{https://www.math.ru.nl/geometryseminar/}{Geometry Seminar},
Radboud University Nijmegen, March 2023.
\item \emph{Moduli of Bridgeland stable triples.} International Centre for
Theoretical Physics, Trieste, March 2022.
\item \emph{Projective moduli of semistable hypersurfaces.}
\href{https://mathematik.univie.ac.at/en/research/seminars/geometry-and-topology-seminar/}{Geometry and Topology Seminar},
Universit\"at Wien, January 2023.
\item \href{https://www.youtube.com/watch?v=5IlL8cVqH2U}{\em Normal simplices and moduli of weighted hypersurfaces.}
\href{https://sites.google.com/view/africa-math-seminar/home}{African Mathematics Seminar}, April 2021.
\item \emph{polymake introduction.}
\href{https://www.mis.mpg.de/calendar/conferences/2020/polymake2020}{11th polymake Conference and Developer Meeting},
MPI for Mathematics in the Sciences, Leipzig, January 2020.
\item \emph{What are algebraic curves and their tropical friends?}
\href{https://whatisseminar.xyz/talks.html}{``What is\ldots?'' Seminar},
Berlin Mathematical School, Urania Berlin, November 2019.
{\em Joint with Marta Panizzut.}
\item \emph{What is $M_{g,n}$?}
\href{https://whatisseminar.xyz/talks.html}{``What is\ldots?'' Seminar},
Berlin Mathematical School, FU Berlin, October 2019.
\item \emph{On the moduli of weighted hypersurfaces.}
Max Planck Institute for Mathematics in the Sciences, Leipzig, 2019.
% TODO: month, if you want it (you thought June or July 2019).
\item \emph{Moduli of hypersurfaces in weighted projective space.}
\href{https://www.maths.ox.ac.uk/node/32016}{Algebraic Geometry Seminar},
University of Oxford, May 2019.
\item \emph{Moduli of hypersurfaces in weighted projective space.}
International Centre for Theoretical Physics, Trieste, January 2019.
\item \emph{On the moduli of weighted hypersurfaces.}
EAUMP--ICTP summer school on homological methods in algebra and geometry,
University of Dar-es-Salaam, July 2018.
\item \emph{Quotients in classification problems.}
American University in Cairo, September 2017.
\end{cvlist}

%% ======================= SUPERVISION =======================
\cvsection{Supervision}
\begin{cvlist}
\item David Schmaling, Bachelor's thesis,
{\em Regular sequences and stabiliser groups of polynomials}, TU Berlin, 2025
\item Josef Cutler, Bachelor's thesis,
{\em Locally nilpotent derivations and $\mathbb{G}_a$-actions}, TU Berlin, 2025
\item Yosime Flood, Bachelor's thesis,
{\em Algebraic curves and the Riemann--Roch theorem}, TU Berlin, 2024
\item Jonathan Wasserman, Master's thesis, {\em The Keel--Mori theorem},
TU Berlin, in progress
\item Aleksandra Kozhevnikova, Master's thesis,
{\em Borel--Moore cohomology of geometric quotients}, TU Berlin, in progress
\item Two Bachelor's theses, TU Berlin, in progress
% TODO: names and titles of the two Bachelor's students once their topics settle.
\end{cvlist}

%% ======================= AFRICA =======================
\cvsection{Mathematics in Africa}
{\small Sustained involvement since 2014, in teaching, organisation, programme
governance and outreach. Individual entries also appear under Teaching,
Organisational work and Research visits.}
\begin{cvlist}
\item Lecturer, Makerere University, Kampala, 2022--23; guest lecturer for a
graduate seminar there, 2023; organiser of a workshop there, 2021, and
co-organiser of a conference there, 2027 (forthcoming)
\item Programme committee, Young African Mathematicians fellowship programme,
MATH+, Berlin; international collaborator, Eastern Africa Algebra Research Group
\item Organiser and mini-course lecturer, workshop
{\em Equivariant methods in geometry}, University of Ibadan, 2024
\item Lecture course on lattice polytopes, {\em Algebra and Geometry from Africa}
(Mathematics for Humanity programme), ICMS Edinburgh, 2025
\item Guest lecturer, AIMS Ghana, 2020 and 2021/22; research visit to AIMS
Rwanda, 2023
\item Facilitator and speaker, EAUMP--ICTP summer school, University of
Dar-es-Salaam, 2018; lecturer, 1st EAALG workshop, 2021
\item Public talks {\em Why mathematics?} in Kampala, 2022; A-level teaching and
university-application mentoring, St.\ Peter's Kubatana High School, Harare,
Zimbabwe, 2014
\end{cvlist}

%% ======================= SERVICE =======================
\cvsection{Service and outreach}
\begin{cvlist}
\item Nominated and prepared the successful application for a MATH+ Hanna
Neumann Fellowship (Caroline Namanya, 2026)
\item Referee for {\em Bulletin of the London Mathematical Society}, {\em Journal
of the London Mathematical Society}, {\em International Mathematics Research
Notices} and {\em Advances in Geometry}
\item Programme committee, Young African Mathematicians (YAM) fellowship
programme, MATH+, Berlin
\item International collaborator,
\href{https://sites.google.com/view/eaalg/international-collaborators}{Eastern Africa Algebra Research Group} (EAALG)
\item \href{https://mathplus.de/news/calculate-with-africa-meridian-podcast-featuring-math-members-bunnett-and-zainelabdeen/}{\em Calculate with Africa: transforming the world through mathematics},
Meridian podcast, Berlin University Alliance, 2024
\item {\em Why mathematics?} Public talks for incoming university students at
Makerere University and for pupils at a state secondary school in Kampala, 2022
% TODO: confirm the year of the two Kampala talks, and the school's name.
% TODO (optional): any committee, editorial, admissions or hiring panel work.
\end{cvlist}

%% ======================= RESEARCH VISITS =======================
\cvsection{Research visits}
\begin{cvlist}
\item Isaac Newton Institute satellite programme
{\em Algebraic Groups, Geometry, Invariants and Related Topics},
University of Essex, invited participant, July 2026
\item Isaac Newton Institute, University of Cambridge, programme
{\em New Equivariant Methods in Algebraic and Differential Geometry},
January -- June 2024
\item AIMS Rwanda, {\em visiting} Jan H\k{a}z\l{}a, January 2023
\item University of Essex, Colchester, {\em visiting} Jes\'us Mart\'inez-Garc\'ia, April 2022
\item International Centre for Theoretical Physics, Trieste,
{\em visiting} Alejandra Rinc\'on, March 2022
\item University of Oxford, {\em visiting} Praise Adeyemo and Bal\'azs Szendr\H{o}i, January 2022
\item International Centre for Theoretical Physics, Trieste,
{\em visiting} Tarig Abdelgadir, January 2019
\end{cvlist}

%% ======================= SOFTWARE =======================
\cvsection{Software}
Available at \url{https://github.com/dombunnett}
\par\vspace{0.12cm}
\begin{cvlist}
\item \href{https://github.com/micjoswig/TropicalModuliData/tree/main/2023-Stacky_fans_and_tropical_moduli_in_polymake}{Stacky fans and tropical moduli in polymake}, \texttt{polymake}\\
{\em Joint with Michael Joswig and Julian Pfeifle.}
\item Reduced submodules of finite dimensional polynomial modules, \texttt{python}
\item Determinantal representations of the cubic discriminant, \texttt{M2}
\item Automorphism group of a projective toric orbifold, \texttt{polymake}
\item Non-normal simplices and weighted projective spaces, \texttt{polymake}
\end{cvlist}

%% ======================= ORGANISATIONAL WORK =======================
\cvsection{Organisational work}
\begin{cvlist}
\item Co-organiser: Conference in algebraic geometry, Makerere University,
Kampala, 2027 (forthcoming)
% TODO: exact conference name and dates once fixed
\item Organiser: Workshop {\em Equivariant methods in geometry},
University of Ibadan, January 2024
\item Organiser: Workshop on homological algebra and discrete geometry,
Makerere University, Kampala, 2021
\item Organiser: Workshop on nonlinear algebra, TU Berlin, 2020
\item Group leader: Cubic surfaces workshop, University of Oslo, 2019
\item Facilitator: EAUMP--ICTP summer school on homological methods in algebra
and geometry, University of Dar-es-Salaam, 2018
\item Organiser: Research seminar --- further topics in geometric invariant
theory, Freie Universit\"at Berlin, Summer Semester 2016
\end{cvlist}

%% ======================= TEACHING =======================
\cvsection{Teaching}
\begin{cvlist}
\item Summer Semester 2026, Algebra II (Commutative Algebra), TU Berlin
\item Summer Semester 2026, Seminar on Linear Algebraic Groups, TU Berlin
\item Winter Semester 2025/26, Algebraic Geometry I, TU Berlin
\item Winter Semester 2025/26, Analysis II for Engineers, TU Berlin
\item Since September 2024, {\em Applied Computational Mathematics},
The Open University (Associate Lecturer)
\item May 2025, Lecture course on lattice polytopes,
{\em Algebra and Geometry from Africa}, ICMS, Edinburgh
\item Summer Semester 2025, Schemes reading course, TU Berlin
\item Summer Semester 2024, Algebra II (Commutative Algebra), TU Berlin
\item Winter Semester 2023/24, Algebra I (Fields and Galois Theory), TU Berlin
\item January 2024, Mini-course on moduli spaces and group actions, workshop
\href{https://sites.google.com/view/bunnett/equivariant-methods-in-geometry-ibadan-2023}{\em Equivariant methods in geometry},
University of Ibadan
\item Winter Semester 2023/24, Algebraic Geometry (remote learning course)
\item August 2023, Graduate seminar on advanced topics in commutative algebra
and geometry, Makerere University
\item Summer Semester 2023, Calculus I and Linear Algebra for Engineers, TU Berlin
\item Summer Semester 2022, Algebraic Geometry, Makerere University
\item Winter Semester 2021/22, Algebra I (Galois Theory, auf Deutsch), TU Berlin
\item Winter Semester 2021/22, Computational introduction to algebraic geometry, AIMS Ghana
\item December 2021, Introduction to homological algebra, EAALG workshop
\item Summer Semester 2020, Derived categories in algebraic geometry, Makerere University
\item 2020, Polyhedral geometry and polymake, Abram Gannibal workshop
\item 2020, Singularities of affine hypersurfaces, Macaulay2 workshop
\item Winter Semester 2017/18, Linear Algebra I (auf Deutsch), FU Berlin
\item Summer Semester 2017, Linear Algebra II (auf Deutsch), FU Berlin
\item Winter Semester 2016/17, Differential Geometry I, Berlin Mathematical School
\item 2014--2015, Undergraduate tutor (analysis, linear algebra, abstract
algebra, calculus), University of Warwick; 2012--2015, private tutor at
A-level and GCSE; Summer 2014, A-level mathematics, St.\ Peter's Kubatana
High School, Harare, Zimbabwe
\end{cvlist}

%% ======================= LANGUAGES =======================
\cvsection{Languages}
English (native) \quad $\cdot$ \quad German (fluent) \quad $\cdot$ \quad Luganda (beginner)

\end{document}
